% Options for packages loaded elsewhere
\PassOptionsToPackage{unicode}{hyperref}
\PassOptionsToPackage{hyphens}{url}
\PassOptionsToPackage{dvipsnames,svgnames,x11names}{xcolor}
%
\documentclass[
  letterpaper,
  DIV=11,
  numbers=noendperiod]{scrartcl}

\usepackage{amsmath,amssymb}
\usepackage{iftex}
\ifPDFTeX
  \usepackage[T1]{fontenc}
  \usepackage[utf8]{inputenc}
  \usepackage{textcomp} % provide euro and other symbols
\else % if luatex or xetex
  \usepackage{unicode-math}
  \defaultfontfeatures{Scale=MatchLowercase}
  \defaultfontfeatures[\rmfamily]{Ligatures=TeX,Scale=1}
\fi
\usepackage{lmodern}
\ifPDFTeX\else  
    % xetex/luatex font selection
\fi
% Use upquote if available, for straight quotes in verbatim environments
\IfFileExists{upquote.sty}{\usepackage{upquote}}{}
\IfFileExists{microtype.sty}{% use microtype if available
  \usepackage[]{microtype}
  \UseMicrotypeSet[protrusion]{basicmath} % disable protrusion for tt fonts
}{}
\makeatletter
\@ifundefined{KOMAClassName}{% if non-KOMA class
  \IfFileExists{parskip.sty}{%
    \usepackage{parskip}
  }{% else
    \setlength{\parindent}{0pt}
    \setlength{\parskip}{6pt plus 2pt minus 1pt}}
}{% if KOMA class
  \KOMAoptions{parskip=half}}
\makeatother
\usepackage{xcolor}
\setlength{\emergencystretch}{3em} % prevent overfull lines
\setcounter{secnumdepth}{-\maxdimen} % remove section numbering
% Make \paragraph and \subparagraph free-standing
\makeatletter
\ifx\paragraph\undefined\else
  \let\oldparagraph\paragraph
  \renewcommand{\paragraph}{
    \@ifstar
      \xxxParagraphStar
      \xxxParagraphNoStar
  }
  \newcommand{\xxxParagraphStar}[1]{\oldparagraph*{#1}\mbox{}}
  \newcommand{\xxxParagraphNoStar}[1]{\oldparagraph{#1}\mbox{}}
\fi
\ifx\subparagraph\undefined\else
  \let\oldsubparagraph\subparagraph
  \renewcommand{\subparagraph}{
    \@ifstar
      \xxxSubParagraphStar
      \xxxSubParagraphNoStar
  }
  \newcommand{\xxxSubParagraphStar}[1]{\oldsubparagraph*{#1}\mbox{}}
  \newcommand{\xxxSubParagraphNoStar}[1]{\oldsubparagraph{#1}\mbox{}}
\fi
\makeatother


\providecommand{\tightlist}{%
  \setlength{\itemsep}{0pt}\setlength{\parskip}{0pt}}\usepackage{longtable,booktabs,array}
\usepackage{calc} % for calculating minipage widths
% Correct order of tables after \paragraph or \subparagraph
\usepackage{etoolbox}
\makeatletter
\patchcmd\longtable{\par}{\if@noskipsec\mbox{}\fi\par}{}{}
\makeatother
% Allow footnotes in longtable head/foot
\IfFileExists{footnotehyper.sty}{\usepackage{footnotehyper}}{\usepackage{footnote}}
\makesavenoteenv{longtable}
\usepackage{graphicx}
\makeatletter
\def\maxwidth{\ifdim\Gin@nat@width>\linewidth\linewidth\else\Gin@nat@width\fi}
\def\maxheight{\ifdim\Gin@nat@height>\textheight\textheight\else\Gin@nat@height\fi}
\makeatother
% Scale images if necessary, so that they will not overflow the page
% margins by default, and it is still possible to overwrite the defaults
% using explicit options in \includegraphics[width, height, ...]{}
\setkeys{Gin}{width=\maxwidth,height=\maxheight,keepaspectratio}
% Set default figure placement to htbp
\makeatletter
\def\fps@figure{htbp}
\makeatother

\KOMAoption{captions}{tableheading}
\makeatletter
\@ifpackageloaded{caption}{}{\usepackage{caption}}
\AtBeginDocument{%
\ifdefined\contentsname
  \renewcommand*\contentsname{Table of contents}
\else
  \newcommand\contentsname{Table of contents}
\fi
\ifdefined\listfigurename
  \renewcommand*\listfigurename{List of Figures}
\else
  \newcommand\listfigurename{List of Figures}
\fi
\ifdefined\listtablename
  \renewcommand*\listtablename{List of Tables}
\else
  \newcommand\listtablename{List of Tables}
\fi
\ifdefined\figurename
  \renewcommand*\figurename{Figure}
\else
  \newcommand\figurename{Figure}
\fi
\ifdefined\tablename
  \renewcommand*\tablename{Table}
\else
  \newcommand\tablename{Table}
\fi
}
\@ifpackageloaded{float}{}{\usepackage{float}}
\floatstyle{ruled}
\@ifundefined{c@chapter}{\newfloat{codelisting}{h}{lop}}{\newfloat{codelisting}{h}{lop}[chapter]}
\floatname{codelisting}{Listing}
\newcommand*\listoflistings{\listof{codelisting}{List of Listings}}
\makeatother
\makeatletter
\makeatother
\makeatletter
\@ifpackageloaded{caption}{}{\usepackage{caption}}
\@ifpackageloaded{subcaption}{}{\usepackage{subcaption}}
\makeatother
\ifLuaTeX
  \usepackage{selnolig}  % disable illegal ligatures
\fi
\usepackage{bookmark}

\IfFileExists{xurl.sty}{\usepackage{xurl}}{} % add URL line breaks if available
\urlstyle{same} % disable monospaced font for URLs
\hypersetup{
  colorlinks=true,
  linkcolor={blue},
  filecolor={Maroon},
  citecolor={Blue},
  urlcolor={Blue},
  pdfcreator={LaTeX via pandoc}}

\author{}
\date{}

\begin{document}

\section{Did Betting Early-Season Underdogs Ever Actually
Work?}\label{did-betting-early-season-underdogs-ever-actually-work}

Spencer's theory is intuitive. At the beginning of a season, the market
still does not know exactly who is good, who is bad, and which preseason
narratives were wrong. If sportsbooks and bettors are still leaning too
hard on stale priors, some teams priced as underdogs should really be
closer to coin flips. And because underdogs pay more when they win, even
a modest amount of mispricing could, in principle, turn into positive
expected return.

That story is plausible. The harder question is whether it was true
often enough, across leagues and through bookmaker vig, to show up in
real historical moneyline data.

I built a cross-league game panel covering the NFL, NBA, NHL, and MLB
from \texttt{2011} through \texttt{2021} using public sportsbook
archives with closing moneylines and final scores. The final cleaned
panel contains \texttt{53,453} games across \texttt{44} league-seasons.
The package, code, and source memo live in the replication folder:
\href{../replication/}{\texttt{../replication/}}.

\subsection{The Strategy Definitions}\label{the-strategy-definitions}

The baseline strategy is deliberately naive:

\begin{itemize}
\tightlist
\item
  bet \texttt{\$100} on every underdog
\item
  use closing moneyline odds
\item
  focus on the very beginning of each season
\end{itemize}

I fixed four early-season definitions before reviewing the output tables
for this project:

\begin{enumerate}
\def\labelenumi{\arabic{enumi}.}
\tightlist
\item
  First 2 calendar weeks.
\item
  First 3 calendar weeks.
\item
  First 5 games played by the underdog.
\item
  First 10 percent of the team-season proxy.
\end{enumerate}

I also compared the rule with obvious benchmarks:

\begin{itemize}
\tightlist
\item
  betting all underdogs over the likely-regular-season proxy
\item
  betting favorites in the same early windows
\item
  betting a random team in each early game
\item
  betting underdogs later in the season
\item
  permuting underdog timing within season to see whether the actual
  early window looks special
\end{itemize}

\subsection{Main Result}\label{main-result}

The best naive pooled rule was ``bet every early underdog for the first
3 weeks.'' It still did not make compelling money.

\begin{figure}[H]

{\centering \includegraphics{./Users/jacobbrown/Documents/GitHub/codex-my-workflow/explorations/spencer-underdog-betting/output/fig_roi_by_window.png}

}

\caption{ROI by window}

\end{figure}%

Pooled results:

\begin{itemize}
\tightlist
\item
  First 2 weeks: ROI \texttt{-1.2\%}
\item
  First 3 weeks: ROI \texttt{-0.4\%}
\item
  First 5 games: ROI \texttt{-1.0\%}
\item
  First 10 percent of season: ROI \texttt{-0.7\%}
\end{itemize}

The first-3-weeks rule placed \texttt{5,901} bets and produced a
\texttt{95\%} bootstrap interval of roughly
\texttt{{[}-4.0\%,\ 3.4\%{]}}. That interval is important. It means the
historical sample is consistent with mild gains, but also with ordinary
losses. The evidence is not strong enough to call this a proven positive
edge.

\subsection{Better Than the Alternatives, But Still Not a Gold
Mine}\label{better-than-the-alternatives-but-still-not-a-gold-mine}

Where Spencer's idea looks strongest is in relative, not absolute,
terms.

Compared with the obvious baselines:

\begin{itemize}
\tightlist
\item
  all underdogs over the likely-regular-season proxy lost \texttt{1.8\%}
\item
  early favorites lost about \texttt{4.1\%}
\item
  a random team in the same early window lost about \texttt{2.4\%} on
  average
\item
  permuted underdog timing also looked worse than the actual early
  window
\end{itemize}

So early underdogs were less bad than generic underdog betting. That
pattern is consistent with a small early-season informational story that
mostly gets eaten by vig and variance, though it is not diagnostic of
that mechanism on its own.

\subsection{The League Split}\label{the-league-split}

The story is not the same in every league.

\begin{figure}[H]

{\centering \includegraphics{./Users/jacobbrown/Documents/GitHub/codex-my-workflow/explorations/spencer-underdog-betting/output/fig_forest_plot.png}

}

\caption{Forest plot}

\end{figure}%

Early 3-week underdog ROI by league:

\begin{itemize}
\tightlist
\item
  MLB: \texttt{+3.9\%}
\item
  NFL: \texttt{+1.8\%}
\item
  NBA: \texttt{-0.4\%}
\item
  NHL: \texttt{-8.2\%}
\end{itemize}

The partial-pooling estimates shrink those league-specific returns back
toward roughly zero, especially for NFL and MLB. That matters because
the raw split is too noisy to interpret as four independent truths.

The cleanest substantive read is:

\begin{itemize}
\tightlist
\item
  MLB shows the most interesting positive hint.
\item
  NFL is directionally positive but very uncertain because the sample is
  small.
\item
  NBA looks close to efficient.
\item
  NHL looks actively hostile to the strategy.
\end{itemize}

That does \textbf{not} line up cleanly with any one simple theory like
``shorter seasons help'' or ``roster continuity helps.'' If there is a
real mechanism, it is more nuanced than one league characteristic.

\subsection{Was the Market Actually Underpricing Early
Underdogs?}\label{was-the-market-actually-underpricing-early-underdogs}

Raw ROI can be misleading because underdogs are paid more when they win.
A better diagnostic is to compare realized upset frequency with the
no-vig implied probability.

For pooled early 3-week underdogs:

\begin{itemize}
\tightlist
\item
  realized win rate: \texttt{40.45\%}
\item
  average no-vig implied win probability: \texttt{39.70\%}
\item
  calibration gap: \texttt{+0.75} percentage points
\end{itemize}

\begin{figure}[H]

{\centering \includegraphics{./Users/jacobbrown/Documents/GitHub/codex-my-workflow/explorations/spencer-underdog-betting/output/fig_calibration.png}

}

\caption{Calibration plot}

\end{figure}%

That is a real but small gap. It says early underdogs won slightly more
often than the de-vigged closing prices implied, but not by enough to
create a large, stable profit margin after transaction costs.

League-level calibration gaps sharpen the picture:

\begin{itemize}
\tightlist
\item
  MLB: \texttt{+2.14} points
\item
  NFL: \texttt{+2.69} points
\item
  NBA: \texttt{+0.65} points
\item
  NHL: \texttt{-2.15} points
\end{itemize}

Again, the same pattern appears: some positive early underdog drift in
MLB and NFL, little in NBA, and negative drift in NHL.

\subsection{A More Refined Version?}\label{a-more-refined-version}

The most interesting refined variant is to focus on bigger underdogs in
the first 3 weeks. That bucket shows a pooled ROI of \texttt{+4.2\%}.

That is the kind of result that can encourage overclaiming. The problem
is that the confidence interval is still wide, the NHL contribution is
strongly negative, and the result appears only after looking across
several reasonable strategy variants. Once you remember multiple-testing
risk, this looks more like a lead than a conclusion.

By contrast, modest underdogs in the first 3 weeks were clearly bad:

\begin{itemize}
\tightlist
\item
  \texttt{modest\_dog\_3w}: ROI \texttt{-3.4\%}
\end{itemize}

So if there is any real edge here, it seems concentrated in the tail of
the underdog distribution rather than in generic ``small dog'' pricing.

\subsection{Kelly Sizing Does Not Rescue the
Story}\label{kelly-sizing-does-not-rescue-the-story}

I also ran a simple out-of-sample half-Kelly exercise using
leave-one-season-out predicted win probabilities. If there were a
meaningful exploitable edge, a careful staking rule should improve the
picture.

It barely does:

\begin{itemize}
\tightlist
\item
  half-Kelly, 2 weeks: \texttt{+0.1\%}
\item
  half-Kelly, 3 weeks: \texttt{+0.2\%}
\end{itemize}

That is basically breakeven. It is another reason to think the true
edge, if it exists at all, is small.

\subsection{A Market-Learning
Simulation}\label{a-market-learning-simulation}

The historical data alone cannot tell us exactly why the strategy
underperforms or occasionally shines. So I added a simple
market-learning simulation with:

\begin{itemize}
\tightlist
\item
  preseason priors about team strength
\item
  slow or fast market updating
\item
  a public premium on favorites
\end{itemize}

In that toy model, early underdogs can become profitable if the market
learns slowly and favorites are overpriced. That mechanism is plausible.
But the historical results are only directionally compatible with that
mechanism; the simulation is too stylized to claim a quantitative match.

So the simulation supports the logic of Spencer's theory, but not a
strong empirical claim that the edge was large in the observed data.

\subsection{Bottom Line}\label{bottom-line}

Spencer's idea was plausible, and the historical record suggests it was
directionally smarter than betting underdogs indiscriminately all year.

But the strongest empirical statement I think the data support is this:

\begin{quote}
Early-season underdogs historically looked slightly less overpriced than
underdogs in general, but not reliably profitable enough to count as a
robust cross-league betting edge.
\end{quote}

If I had to summarize the verdict:

\begin{itemize}
\tightlist
\item
  Did naive early underdogs make money historically? Probably not in a
  robust cross-league sense.
\item
  Was the strategy better than generic underdog betting? Yes.
\item
  Is there a weak hint of early-season underdog underpricing in MLB and
  maybe NFL? Also yes.
\item
  Should a serious bettor trust this as a durable edge? Not on this
  evidence.
\end{itemize}

\subsection{Replication}\label{replication}

The code-only replication package for this post lives here:

\begin{itemize}
\tightlist
\item
  \href{../replication/}{\texttt{../replication/}}
\end{itemize}

It contains:

\begin{itemize}
\tightlist
\item
  download scripts for the public raw data
\item
  the analysis pipeline
\item
  the R robustness script
\item
  dependency declarations
\item
  a one-command rebuild script
\end{itemize}



\end{document}
